\documentclass[10pt,a4paper]{amsart}
\usepackage[margin=19mm]{geometry}
\usepackage{amsmath,amssymb,mathtools}
\usepackage[scr=rsfs,cal=euler]{mathalfa}
\usepackage{xfrac}
\usepackage[unicode,hidelinks,bookmarks=false]{hyperref}
\newcommand{\ie}{i.e.\ }
\newcommand{\cf}{cf.\ }
\newcommand{\resp}{respectively\ }
\newcommand{\Subsection}{\subsection}
\providecommand{\nobreakdash}{\nobreak\hbox{-}}
\setlength{\emergencystretch}{2em}
\multlinegap0pt
\usepackage{amssymb}
\usepackage{mathtools}
\usepackage{enumerate}
\usepackage{tikz}
\usepackage{tcolorbox}
\newtheorem{proposition}{Proposition}
\DeclareMathOperator{\FBLp}{FBL^{\mathit{(p)}}}
\newcommand{\R}{\mathbb{R}}
\title{Upper bound properties of free and related Banach lattices via operators}
\author{Anastasiia Ianina, Timur Oikhberg}
\date{Source: arXiv:2609.11928v1 (CC BY 4.0)}
\begin{document}
\maketitle

\noindent\emph{Source and licence.} Upper bound properties of free and related Banach lattices via operators. Version arXiv:2609.11928v1 (CC BY 4.0), distributed by arXiv under the Creative Commons Attribution 4.0 International licence, http://creativecommons.org/licenses/by/4.0/, as recorded in the arXiv licence metadata for this exact version and shown on the abstract page https://arxiv.org/abs/2609.11928v1. Full licence notice: \texttt{licenses/arxiv-2609.11928-CC-BY-4.0.txt}.\par\smallskip

\noindent\textit{Editorial scope.} Counted unit: the article's proposition FBL (source0.tex lines 583--585). The article's complete original proof of it is delivered with it (source0.tex lines 587--629, one \texttt{proof} environment of 43 lines). No mathematics is written, repaired or re-derived: the statement and the proof are the article's own lines, in the article's own notation, and no retained line is substituted or re-flowed.  No further source-local context is needed: every cross-reference of every retained line resolves inside the two delivered spans.  Nothing was omitted from any retained span, so the package carries no omission record.  No retained line contains a citation, so the wrapper carries no bibliography and no bibliography entry.  No retained line contains a figure environment, an image-inclusion command or drawing code, so no image is delivered and the package declares no asset dependency.  Editorial formatting only, in addition: an AMS \texttt{amsart} preamble that restates the article's own declarations and macros used by the retained lines, namely the environments \texttt{proposition} on separate counters, together with the standard packages those lines need.  The retained environments print in this document as Proposition 1 in order of appearance, whereas the article prints the same items with the numbering of its own full document.  The complete native source of the article as distributed in the arXiv e-print for this exact version is bundled unchanged as \texttt{sources/210062/source0.tex}.
\begin{proposition}\label{supremum in FBL}
Suppose $T \in B(Z,X)$ is compact and dual $p$-persistent relative to a dense $X_0 \subset X$. Define $\phi : X^* \to \R : x^* \mapsto \|T^*x^*\|$. Fix a norm one $x_0 \in X_0$ and $\mu > 0$. If $g \in \FBLp[X]$ satisfies $g \geq f = \mu \big| \delta_{x_0} \big| \wedge \phi$, then $g \geq \mu \big| \delta_{x_0} \big|$.
\end{proposition}

\begin{proof}
By scaling, we assume that $c$ in the definition of dual $p$-persistence equals $1$.
Pick a norm one $x^* \in X^*$, and show that $g(x^*) \geq \mu |\langle x^*, x_0\rangle|$. 

Pick $\delta > 0$, and write $g = g_1 + g_2$, where $g_1 = G(\delta_{x_1}, \ldots, \delta_{x_m})$, $G$ being a lattice polynomial, and $\|g_2\| < \delta$.
Find $x_1^*, \ldots, x_n^* \in \cap_{i=0}^m x_i^\perp$ so that $\|(x_j^*)\|_{p,{\textrm{weak}}} \leq 1$, and $\sum_j \alpha_j^p = 1$, where $\alpha_j = \|T^* x_j^*\|$.

Pick $C > \|T\| + \mu + 1$. For each $j$, $g(\alpha_j x^* + C x_j^*) = \alpha_j g_1(x^*) + g_2(\alpha_j x^* + C x_j^*)$, or, in other words,
$$
\big( g(\alpha_j x^* + C x_j^*) \big)_j - \big( \alpha_j \big)_j g_1(x^*) = \big( g_2(\alpha_j x^* + C x_j^*) \big)_j .
$$
The triangle inequality in $\ell_p^n$ gives us:
$$
\Big|\big\| \big( g(\alpha_j x^* + C x_j^*) \big)_j \big\|_p - g_1(x^*) \Big| \leq \big\| \big( g_2(\alpha_j x^* + C x_j^*) \big)_j \big\|_p 
$$
(here we use our scaling: $\big\| \big( \alpha_j \big)_j \big\|_p = 1$). Taking $\|x^*\| = 1$ into account,
$$
\big\| \big( \alpha_j x^* + C x_j^* \big)_j \big\|_{p, {\textrm{weak}}} \leq C \big\| \big( x_j^* \big)_j \big\|_{p, {\textrm{weak}}} + \big\| \big( \alpha_j \big)_j \big\|_p \leq C+1 ,
$$
hence $\big\| \big( g_2(\alpha_j x^* + C x_j^*) \big)_j \big\|_p < \delta (C+1)$, and therefore,
\begin{equation}
\Big|\big\| \big( g(\alpha_j x^* + C x_j^*) \big)_j \big\|_p - g_1(x^*) \Big| < (C+1)\delta .
\label{g g1 close}
\end{equation}

Now note that
$$
f(\alpha_j x^* + C x_j^*) = \mu \alpha_j \big| \langle x^*, x_0 \rangle \big| \wedge \|T^* (\alpha_j x^* + C x_j^*)\| ,
$$
and
$$
\|T^* (\alpha_j x^* + C x_j^*)\| \geq C \|T^* x_j^*\| - \alpha_j \|T^* x^*\| = \alpha_j (C - \|T^*x^*\|) .
$$
By our choice of $C$, $C - \|T^*x^*\| > \mu + 1$, hence $f(\alpha_j x^* + C x_j^*) = \mu \alpha_j \big| \langle x^*, x_0 \rangle \big|$. We clearly have
$$
\big\| \big( g(\alpha_j x^* + C x_j^*) \big)_j \big\|_p \geq \big\| \big( f(\alpha_j x^* + C x_j^*) \big)_j \big\|_p = \mu \big\| \big( \alpha_j \big)_j \big\|_p \big| \langle x^*, x_0 \rangle \big| = \mu \big| \langle x^*, x_0 \rangle \big| ,
$$
so from \eqref{g g1 close}, $g_1(x^*) \geq \mu \big| \langle x^*, x_0 \rangle \big| - (C+1)\delta$. Hence,
$$
g(x^*) \geq g_1(x^*) - \|g_2\| \geq \mu \big| \langle x^*, x_0 \rangle \big| - (C+2)\delta .
$$
To complete the proof, recall that $\delta$ can be arbitrarily small.
\end{proof}

\end{document}
